\documentclass[11pt]{article}
\usepackage[T1]{fontenc}
\usepackage[utf8]{inputenc}
\usepackage[margin=0.9in]{geometry}
\usepackage{amsmath,amssymb,amsthm,microtype}
\usepackage[hidelinks]{hyperref}

\newtheorem{theorem}{Theorem}
\newtheorem{lemma}[theorem]{Lemma}
\newtheorem{proposition}[theorem]{Proposition}
\newtheorem{corollary}[theorem]{Corollary}

\title{Every two-vertex premaniplex is the symmetry type graph of a finite abstract polytope}
\author{Carptopus\\\small\texttt{carptopus@163.com}}
\date{Public Beta, 3 September 2026}

\begin{document}
\maketitle
\begin{abstract}

For every rank $r\geq 3$, there are $2^r-1$ connected two-vertex premaniplexes, one for each
nonempty set of colors assigned to links between the two vertices. Pellicer, Poto\v{c}nik and Toledo
proved that every such premaniplex occurs as the symmetry type graph of a finite maniplex. The
corresponding realization problem for finite abstract polytopes remained open. Mochán established
polytopal realizations when the link-color set has size one or two, or is an interval, and reduced
the remaining cases in his voltage construction to one family of path-intersection identities.

We prove that every connected two-vertex premaniplex is the full symmetry type graph of a finite
abstract polytope. The main step is a support-localization argument for the recursive regular
polytopes used by Mochán. It treats an arbitrary set of internal link colors simultaneously and
establishes the missing $k=1$ path-intersection identity. The two exceptional rank-four base types
are handled by exact finite computation. Thus the two-vertex realization problem has an affirmative
answer in every rank at least three.

\end{abstract}

\section{Introduction}

An abstract polytope is a ranked, strongly flag-connected partially ordered set satisfying the
diamond condition. Its flags form a properly edge-colored graph: two flags are joined by an edge of
color $i$ when they differ only in their rank-$i$ face. The quotient of this flag graph by the full
automorphism group is the symmetry type graph of the polytope. Its vertices are the flag orbits,
and a color is represented either by a link between two orbits or by a semi-edge at one orbit.

Premaniplexes axiomatize the colored quotient graphs that may occur in this way. For a rank-$r$
premaniplex with exactly two vertices, every color $i\in\{0,\ldots,r-1\}$ is either a link between
the two vertices or a semi-edge at both vertices. The graph is connected precisely when at least
one color is a link. Hence there are $2^r-1$ connected two-vertex candidates.

Pellicer, Poto\v{c}nik and Toledo \cite{ppt2019} proved that all these candidates are realized by finite
maniplexes. Polytopality is stronger: a maniplex is the flag graph of an abstract polytope exactly
when it satisfies the path-intersection property of Garza-Vargas and Hubard \cite{gvh2018}. Hubard and Mochán
\cite{hm2023} translated this property into an intersection condition for voltage sets over a premaniplex.

Mochán \cite{mochan2024} applied that framework to the two-vertex problem. He constructed finite polytopes when
the link-color set has one or two elements or is an interval. For a general two-vertex type whose
extreme link colors have been normalized to $0$ and $n$, his Theorem 6.4 proves every required
voltage intersection with $k>1$. The unresolved part is the $k=1$ identity. The difficulty noted in
\cite{mochan2024} is that an internal link of color $p$ has voltage involving $r_0r_p$, so a path that avoids color
$0$ may nevertheless move $0$-faces.

The present argument controls exactly that movement. After deleting color $0$, the upper voltage
group is generated by local semi-edge voltages and one fundamental-cycle voltage for every internal
link. All these generators preserve a pair of endpoint regions in the recursive $\widehat 2$
construction. Commutation with regular automorphisms then converts this coarse endpoint control
into the exact path-intersection identity.

Our main result is the following.

\begin{theorem}

For every $r\geq3$, every connected rank-$r$ two-vertex premaniplex is the symmetry type graph,
with respect to the full automorphism group, of a finite abstract polytope.

The proof in ranks at least five is structural. Two rank-four types not covered by the previously
known families are finite computer-assisted base cases. The computation is described in Section 7,
including its reproducibility and independence limitations.

\end{theorem}

\section{Premaniplexes, voltages and the intersection criterion}

We use right actions on flags. If $M$ is a regular rank-$r$ polytope, write

\[
\mathrm{Mon}(M)=\langle r_0,\ldots,r_{r-1}\rangle
\]

for its monodromy group, where $r_i$ sends a flag to its $i$-adjacent flag. Write $\rho_i$ for the
corresponding distinguished automorphisms. Monodromies commute with automorphisms on the flag set.

Let $X$ be a premaniplex and let $\Pi_I^{s,t}(X)$ denote the set of maniplex-homotopy classes of
paths from the vertex $s$ to the vertex $t$ using only colors in $I$. A voltage assignment is an
antimorphism

\[
\xi:\Pi(X)\longrightarrow G.
\]

Thus, if paths are traversed from left to right, their dart voltages multiply in reverse order.
The derived colored graph $X^\xi$ has vertices $(s,g)$ and lifts each dart according to its voltage.

For the voltage assignments used below, the local simplicity and distant-color square conditions
are already part of Mochán's construction. The remaining polytopality criterion is the following
interval form of the weak path-intersection property \cite[Theorem 4.2]{hm2023}: for all vertices $s,t$ and all
$0\leq k,m\leq n$,

\[
\xi\bigl(\Pi_{[0,m]}^{s,t}(X)\bigr)
\cap
\xi\bigl(\Pi_{[k,n]}^{s,t}(X)\bigr)
=
\xi\bigl(\Pi_{[k,m]}^{s,t}(X)\bigr).
\tag{1}
\]

Here $[k,m]$ is empty when $k>m$; the voltage set of the empty color set is the identity when
$s=t$ and is empty otherwise.

For a two-vertex premaniplex $X$, denote its two vertices by $a$ and $b$, and denote the set of link
colors by $L(X)$. Every other color is a semi-edge at both vertices. We shall first solve the
normalized core

\[
L(X)=\{0,n\}\cup P,
\qquad P\subseteq\{1,\ldots,n-1\},
\tag{2}
\]

where $n\geq4$. Lower-rank cores and colors outside the extreme links are treated in Section 8.


\section{The recursive regular-polytope construction}

We recall only the parts of Mochán's construction \cite[Sections 4--6]{mochan2024} that enter the new argument.
Starting from a square, recursively apply the $\widehat 2$ operation to obtain finite regular
polytopes $M_n$. A flag of $M_n=\widehat 2M_{n-1}$ can be written as

\[
(\phi,x),\qquad
x\in\mathbb F_2^{\mathrm{Fac}(M_{n-1})}.
\]

The support $\mathrm{supp}(x)$ is the set of facets of $M_{n-1}$ on which $x$ is nonzero. The
construction is a lattice, and equality of lower faces is governed by containment of this support
in facet stars. Mochán defines a distinguished involutory flag-permutation
$T=\widetilde\rho_0$. Its restriction to each facet is the facet automorphism reflecting across the
chosen base edge; $T$ need not be a global automorphism of $M_n$. Corollary 5.2 gives the local
feature used below: whenever the support is contained in the union of the facet stars of two faces,
the current facet has the original base edge and $T$ acts there as the prescribed reflection.

Let $X_n$ be obtained from $X$ by deleting its color-$n$ link. Mochán's construction starts with a
covering

\[
p:M_n\longrightarrow X_n.
\]

Color the flags of $M_n$ white or black according as their images under $p$ are $a$ or $b$, and
choose the base flag $\Phi$ to be white. Every monodromy of $M_n$ either preserves both color
classes or interchanges them. The voltages below preserve the two classes. In particular, projecting
a flag path in $M_n$ through $p$ records both its color word and whether its endpoints project to
the same or different vertices of $X$.

Let $\Phi=(\phi,0)$ be the base flag of $M_n$. Put

\[
q_a=1,\qquad q_b=r_0,\qquad y=Tr_0.
\tag{3}
\]

The involutions $T$ and $r_0$ commute, so the order in the last product is immaterial. The color-$n$
link has voltage $y$. For every dart $d$ of color $c<n$, directed from $s$ to $t$, Mochán's voltage
assignment satisfies

\[
\xi(d)=q_t r_c q_s.
\tag{4}
\]

This includes the color-$0$ link, whose two directed voltages are both trivial. If a path $W$ from
$s$ to $t$ has color word $c_1\cdots c_\ell$ and

\[
h=r_{c_1}\cdots r_{c_\ell},
\]

then telescoping (4), with the antimorphism convention, gives

\[
\xi(W)=q_t h^{-1}q_s.
\tag{5}
\]

Mochán also augments $\xi$ by a finite cyclic coordinate to obtain $\xi'$. This auxiliary coordinate
prevents an unwanted automorphism from interchanging the two fibers, so that the full symmetry type
graph of the derived maniplex is exactly $X$
\cite[equation (3) and Corollary 6.5]{mochan2024};
\cite[Propositions 19 and 20 and Theorems 29 and 31]{ppt2019}.
For $m<n$, paths using colors in $[0,m]$ do not change the auxiliary coordinate. Consequently, once
(1) is proved for $\xi$, the projection argument of \cite[Corollary 6.5]{mochan2024} proves it for $\xi'$ as well.

Finally, \cite[Theorem 6.4]{mochan2024} proves (1) for every $k>1$, every $m$, every pair of endpoints and every
link set of the form (2). Thus only $k=1$ remains.


\section{The upper voltage group for arbitrary internal links}

Delete color $0$ from $X$ and choose the color-$n$ link as a spanning tree. Let

\[
H_P=\xi\bigl(\Pi_{[1,n]}^{a,a}(X)\bigr)
\]

be the closed upper voltage group at $a$.

\begin{lemma}

The group $H_P$ is generated by

\[
\begin{aligned}
r_i,\;Tr_iT &&& (i\in\{1,\ldots,n-1\}\setminus P),\\
Tr_p &&& (p\in P).
\end{aligned}
\tag{6}
\]

\end{lemma}

\begin{proof}

A semi-edge of color $i$ at $a$ contributes $r_i$. The corresponding semi-edge at $b$, transported
back along the spanning-tree link, contributes $Tr_iT$. For an internal link of color $p$, the
fundamental circuit formed with the spanning-tree link has voltage

\[
y(r_0r_p)=Tr_p.
\]

A connected two-vertex graph with a fixed spanning tree has fundamental group generated by all its
semi-edges and by one fundamental circuit for every non-tree link. Therefore the displayed list is
complete. In particular, multiple internal links introduce no additional mixed generator.
\end{proof}

The importance of Lemma 2 is that every generator admits the same local support control, regardless
of the size or shape of $P$.


\section{Endpoint-support localization}

Let $e$ be the base edge of $M_{n-1}$, with endpoints $u$ and $v$. Let the base two-face be the
square with cyclically ordered vertices

\[
u,v,w,q.
\]

For a vertex $z$ of $M_{n-1}$, write $\mathcal F(z)$ for the set of facets incident with $z$.

\begin{lemma}[localization]

For every $h\in H_P$:

\begin{enumerate}
\item the first-coordinate vertices of $\Phi h$ and $\Phi^1h$ lie in $\{u,v\}$;
\item the first-coordinate vertex of $\Phi^{010}h$ lies in $\{w,q\}$;
\item in each case the vector-coordinate difference is supported in the union of the facet stars of
the indicated pair.
\end{enumerate}

\end{lemma}

\begin{proof}

We induct on a word in the generators (6) and their inverses. For $i<n-1$, the monodromy $r_i$
does not change the first-coordinate vertex or enlarge its support outside the current endpoint
pair. The monodromy $r_{n-1}$ changes one vector coordinate belonging to a facet through the current
vertex and therefore remains inside the corresponding facet-star union.

Whenever $T$ occurs, the induction hypothesis places the current support in the union of the two
facet stars adjacent to the base edge. Mochán's Corollary 5.2 then says that the current facet still
has $e$ as its base edge. Hence $T$ exchanges $u$ with $v$ and exchanges $w$ with $q$, preserving
the asserted regions.

This proves the claim for $r_i$, $Tr_iT$ and $Tr_p$. It also covers arbitrary interleavings. For
$p<n-1$, $T$ commutes with $r_p$, so the inverse of $Tr_p$ has the same form. At the boundary
$p=n-1$ one must also check

\[
(Tr_{n-1})^{-1}=r_{n-1}T.
\]

Here $r_{n-1}$ first changes the support inside the facet star of the current endpoint; Corollary 5.2
then applies before $T$ acts. Thus the inverse preserves the same invariant.
\end{proof}

The localization lemma distinguishes which endpoint a voltage can reach once the voltage is also a
monodromy.

\begin{lemma}[endpoint classification]

Let $h\in H_P$.

\begin{enumerate}
\item If $h$ is a monodromy of $M_n$, then the first-coordinate vertex of $\Phi h$ is $u$.
\item If $yh$ is a monodromy of $M_n$, then the first-coordinate vertex of $\Phi yh$ is $v$.
\end{enumerate}

\end{lemma}

\begin{proof}

Suppose first that $h$ is a monodromy and that $\Phi h$ has first-coordinate vertex $v$.
Monodromies commute with the automorphism $\rho_1$, so

\[
\Phi^1h=(\Phi h)\rho_1.
\]

The automorphism $\rho_1$ fixes $u,w$ and exchanges $v,q$ on the base square. Thus the right-hand
side has first-coordinate vertex $q$, contradicting Lemma 3, which places the vertex of $\Phi^1h$
in $\{u,v\}$.

For the second assertion, the construction gives

\[
\Phi y=\Phi,\qquad \Phi^1y=\Phi^{010}.
\tag{7}
\]

If $yh$ were a monodromy taking the first-coordinate vertex of $\Phi$ to $u$, commutation with
$\rho_1$ would force $\Phi^1yh$ to have first-coordinate vertex $u$. But by (7) and Lemma 3 the
latter vertex lies in $\{w,q\}$. Hence the first-coordinate vertex of $\Phi yh$ must be $v$.

\end{proof}


\section{The missing \texorpdfstring{$k=1$}{k=1} intersection}

\begin{proposition}

Let $n\geq4$ and let $X$ be the normalized two-vertex premaniplex (2). For every
$0\leq m\leq n$ and every $s,t\in\{a,b\}$,

\[
\xi\bigl(\Pi_{[0,m]}^{s,t}(X)\bigr)
\cap
\xi\bigl(\Pi_{[1,n]}^{s,t}(X)\bigr)
=
\xi\bigl(\Pi_{[1,m]}^{s,t}(X)\bigr).
\tag{8}
\]

\end{proposition}

\begin{proof}

The case $m=n$ is tautological. Assume $m<n$ and begin at $s=a$. Let $\omega$ belong to the
left-hand side of (8). Since $\omega$ is the voltage of a lower path using colors in $[0,m]$, formula
(5) identifies it with a monodromy in $\langle r_0,\ldots,r_m\rangle$. In particular, it fixes all
faces of ranks greater than $m$.

An upper voltage from $a$ to $a$ has the form $h\in H_P$, whereas one from $a$ to $b$ has the form
$yh$ with $h\in H_P$. By Lemma 4, if

\[
\Phi\omega=(\psi,x),
\]

then its first-coordinate vertex is $u$ in the closed case and $v$ in the open case.

We next upgrade this first-coordinate statement to equality of the complete rank-$0$ face in
$M_n=\widehat2M_{n-1}$. If $m<n-1$, the lower monodromy fixes the facet, hence $x=0$. Suppose
$m=n-1$ and abbreviate

\[
U=\mathcal F(u),\quad V=\mathcal F(v),\quad
W=\mathcal F(w),\quad Q=\mathcal F(q).
\]

In the closed case, direct localization gives

\[
\mathrm{supp}(x)\subseteq U\cup V.
\]

Apply the same localization to $\Phi^1\omega=(\Phi\omega)\rho_1$ and pull the support back by
$\rho_1$. This gives

\[
\mathrm{supp}(x)\subseteq U\cup Q.
\]

Every facet containing both $v$ and $q$ contains their join, namely the base two-face, and hence also
contains $u$. Therefore

\[
(U\cup V)\cap(U\cup Q)=U.
\]

Thus $\mathrm{supp}(x)\subseteq U$, which means $(u,x)=(u,0)$ as rank-$0$ faces of the
$\widehat2$ construction.

In the open case, use (7). The direct test gives support in $U\cup V$, while the test from
$\Phi^1y=\Phi^{010}$, followed by the $\rho_1$ pullback, gives support in $V\cup W$. Every facet
containing $u$ and $w$ contains the base two-face and hence $v$, so

\[
(U\cup V)\cap(V\cup W)=V.
\]

Consequently $(v,x)=(v,0)$ as rank-$0$ faces.

In the closed case, $\Phi$ and $\Phi\omega$ now have the same rank-$0$ face and the same faces of
all ranks greater than $m$. Strong flag-connectivity of the intervening section supplies a flag path
$R$ from $\Phi$ to $\Phi\omega$ using only colors $[1,m]$. Let $g$ be its monodromy word and
let $W=p(R)$. Then $W$ is an $a$-to-$a$ path and (5) gives $\xi(W)=g^{-1}$. Since
$\Phi g=\Phi\omega$ and $M_n$ is regular, $g=\omega$; hence the reversed projected path $W^{-1}$
has voltage $\omega$.

In the open case, $\Phi^0=\Phi r_0$ and $\Phi\omega$ have the same rank-$0$ face and the same faces
above rank $m$. Let $R$ be the resulting $[1,m]$ path from $\Phi^0$ to $\Phi\omega$, with monodromy
word $g$. Its projection $W=p(R)$ runs from $b$ to $a$, and (5) gives
$\xi(W)=g^{-1}r_0$. Therefore $W^{-1}$ runs from $a$ to $b$ and has voltage $r_0g$.
The equality $\Phi r_0g=\Phi\omega$, together with regularity of $M_n$, implies $r_0g=\omega$.
Thus $W^{-1}$ is the required path. If $P\cap[1,m]$ is empty, no $[1,m]$ path can switch the two
vertices of $X$. The supposed intersection element would therefore produce an impossible projected
path, so both sides of (8) are empty in this boundary case.

This proves (8) for paths starting at $a$. The color-preserving automorphism $\sigma$ of $X$ that
interchanges $a$ and $b$ satisfies, for every dart $d$,

\[
\xi(\sigma d)=r_0\xi(d)r_0.
\tag{9}
\]

For the top color this follows because $T$ commutes with $r_0$. Conjugating the three voltage sets in
(8) by $r_0$ proves the two remaining ordered endpoint cases. The reverse inclusion in (8) is
automatic, completing the proof.
\end{proof}

\begin{corollary}

For every normalized core (2), both $\xi$ and the augmented assignment $\xi'$ satisfy the full
intersection condition (1).

\end{corollary}

\begin{proof}

Proposition 5 proves $k=1$. The case $k=0$ and the case $m=n$ are tautological. Every case $k>1$
is \cite[Theorem 6.4]{mochan2024}. This proves (1) for $\xi$.

If $m<n$, neither a lower path nor the reconstructed $[k,m]$ path uses the top color, so both act
trivially on the auxiliary cyclic coordinate of $\xi'$. The projection argument of
\cite[Corollary 6.5]{mochan2024} therefore lifts (1) from $\xi$ to $\xi'$. The case $m=n$ remains
tautological.
\end{proof}

By \cite[Theorem 4.2]{hm2023}, the derived maniplex $X^{\xi'}$ is polytopal. The finite cyclic augmentation,
\cite[Propositions 19 and 20]{ppt2019} and \cite[Theorems 29 and 31]{ppt2019} ensure, for a sufficiently large finite
parameter, that no extra automorphism merges the two fibers. Hence its full symmetry type graph is
exactly $X$.


\section{The rank-four base cases}

When the distance between the extreme link colors is $n=3$, the only normalized cases not already
covered by one or two links or by an interval are

\[
L(X)=\{0,1,3\},\qquad L(X)=\{0,2,3\}.
\tag{10}
\]

We verified these cases using the regular toroidal polyhedron $\{4,4\}_{(8,0)}$, with 512 flags, as
the rank-three input. The separating involutions are represented by the color words

\[
2101210212\quad\text{and}\quad210212,
\]

respectively. For both types the verification checks:

\begin{enumerate}
\item the defining distant-color square relations and the original base-flag conditions;
\item that every semi-edge voltage is a nonidentity involution;
\item that parallel link colors have distinct voltages at each endpoint;
\item every nontrivial pair $k=1,2,3$ and $m=0,1,2$ in (1);
\item all four ordered endpoint pairs;
\item the known two-link type $\{0,3\}$ as a positive control;
\item failure after replacing the top voltage by the identity as a destructive negative control.
\end{enumerate}

Two implementations reconstruct the flags, reflections, base flags and voltage darts independently.
They do not import one another. For the largest upper group, of order $2^{21}$, both use SymPy's
deterministic Schreier--Sims implementation for membership testing; their group-theoretic kernels
are therefore not independent. Every explicit group closure is preceded by an exact order
calculation and refused above 4096; every remaining reachable-state enumeration has the same
4096-item hard cap. The supplied runner executes the programs serially, monitors the complete
process tree, and terminates it if sampled private memory exceeds 1024 MiB or the run exceeds
120 seconds.

The resource guard itself has a destructive test: a 1 MiB limit terminates the process tree and
leaves no Python process behind. Under the standard guard, complete reproduction used 64.8 MiB and
58.1 MiB of peak private memory for the two implementations, with running times 10.47 seconds and
4.90 seconds. Both exited with code zero, every exact intersection comparison passed, and the
negative control failed as required.

This section is a computer-assisted finite lemma. It does not replace the general proof in Sections
4--6, and it depends on the correctness of the two implementations and of their shared SymPy
membership kernel.


\section{Proof of the main theorem}

\begin{proof}

Let $X$ be a connected rank-$r$ two-vertex premaniplex and let $L=L(X)$ be its nonempty link-color
set.

All rank-three types are known to be realized by finite abstract polyhedra. If $|L|\leq2$, or if
$L$ is an interval, the claim is \cite[Theorem 6.9]{mochan2024}. Otherwise let $i=\min L$ and $j=\max L$.
Restrict to the colors $[i,j]$ and shift their labels so that the extreme links become $0$ and
$n=j-i$.

If $n\leq2$, the core is among the already known cases. If $n=3$, the only remaining cores are
(10), which are polytopal by Section 7. If $n\geq4$, Corollary 6 and the voltage criterion produce a
finite abstract polytope whose full symmetry type graph is the normalized core.

Every color below $i$ or above $j$ is necessarily a semi-edge color. Repeated application of the
standard $\widehat2$ rank-raising operation and its dual restores these exterior semi-edge colors.
These operations preserve finiteness and the number of flag orbits and add precisely the required
semi-edge to the symmetry type graph \cite[Section 4]{mochan2024}; \cite[Proposition 11]{ppt2019}. Thus the resulting finite
abstract polytope has full symmetry type graph $X$.
\end{proof}


\section{Scope and provenance}

The contribution claimed here is Theorem 1, specifically the arbitrary-internal-link localization
in Lemmas 2--4 and the resulting $k=1$ intersection proof in Proposition 5. The voltage criterion,
the recursive regular polytopes $M_n$, the assignments $\xi$ and $\xi'$, the $k>1$ intersection
theorem, and the automorphism-suppression augmentation are cited ingredients from \cite{hm2023,mochan2024,ppt2019}.

Theorem 1 strengthens the maniplex existence theorem of \cite{ppt2019} to abstract polytopes and completes the
two-vertex cases left open after \cite{mochan2024}. The structural theory of two-orbit polytopes developed by
Hubard and Schulte \cite{hs2026} is complementary: it studies the internal structure and automorphism groups
of two-orbit polytopes, rather than supplying this realization theorem.

The result concerns symmetry type graphs with exactly two vertices. It does not classify the
polytopes realizing a given type, bound their minimum size, or address arbitrary symmetry type
graphs with three or more vertices.


\begin{thebibliography}{9}

\bibitem{gvh2018}
J. Garza-Vargas and I. Hubard,
\emph{Polytopality of maniplexes},
Discrete Mathematics 341 (2018), 2068--2079.

\bibitem{hm2023}
I. Hubard and E. Moch\'an,
\emph{All polytopes are coset geometries: characterizing automorphism groups of
$k$-orbit abstract polytopes},
European Journal of Combinatorics 113 (2023), 103746;
arXiv:2208.00547.

\bibitem{mochan2024}
E. Moch\'an,
\emph{Polytopality of 2-orbit maniplexes},
Discrete Mathematics 347 (2024), 114055;
arXiv:2309.15791v3.

\bibitem{ppt2019}
D. Pellicer, P. Poto\v{c}nik and M. Toledo,
\emph{An existence result on two-orbit maniplexes},
arXiv:1812.04148.

\bibitem{hs2026}
I. Hubard and E. Schulte,
\emph{Two-Orbit Polytopes},
arXiv:2604.00185v1, 2026.

\end{thebibliography}

\section*{AI-assisted research and writing disclosure}

Carptopus is the author responsible for the mathematical claims and the public version of this
manuscript. Codex was used as an AI-assisted research, verification and writing tool.

\end{document}
